\begin{abstract}
CASIA~v2.0 is one of the most widely used benchmarks for image forgery detection, and error-level-analysis (ELA)
classifiers commonly report accuracies around 90\% on it. We show that the benchmark as released can be solved
without looking at image content: file metadata alone separates authentic from tampered images with 0.990
balanced accuracy (AUC 0.999), because most tampered images are uncompressed TIFF files while the authentic images
are JPEGs. We therefore evaluate under leak-controlled protocols that re-encode every image (B) or resample and then
re-encode it (R), and we measure every method against the shortcut features that survive each protocol. Our
detector is built from eight interpretable classical techniques covering point processing, histogram processing,
spatial filtering, frequency-domain analysis and edge/keypoint detection: ELA, local histograms, noise, JPEG
ghosts, DCT double quantisation, edge sharpness and two copy-move matchers. Their evidence is summarised by
within-image inconsistency statistics and fused by a random forest. On a frozen test set evaluated once, it reaches
AUC 0.805~[0.779,~0.831] under the strict protocol R, against 0.712 for the shortcut features. The calibrated
detector judges 52\% of images with 86\% accuracy, and learned map fusion localises tampered regions with pixel
F1~0.198 (0.308 under B). The detector transfers to the unseen MICC-F220 copy-move dataset (AUC 0.972; 11 forged scenes). An ELA +
ResNet-18 baseline reaches AUC 0.994 only on the leaking files, where shortcut features alone reach 1.000;
under R it reaches 0.770, and on MICC-F220 its AUC (0.53) is consistent with chance. The system is released with an interactive interface that explains
every verdict.
\end{abstract}

\begin{IEEEkeywords}
image forensics, forgery detection, copy-move, splicing, error level analysis, dataset bias, CASIA v2.0
\end{IEEEkeywords}
